% answers.tex -- marking criteria and sample answers for questions.tex (Module 8 short answer).
% NESA marking-guideline style: a criteria table for each part, then a sample answer.
% Every number: checks/m8-short-answer/numbers.py
% Build: TEXINPUTS=.: xelatex answers.tex (twice)
\documentclass[12pt]{article}
\usepackage{mathtools}
\usepackage{quiz}
\hypersetup{pdftitle={Short Answer Questions - Module 8 - From the Universe to the Atom - Marking guidelines}}
% links between the two PDF files are not possible, so question headings are plain text here
\renewcommand\hyperlink[2]{#2}
\newcommand\anstopic[1]{\par\Needspace{8\baselineskip}\vspace{4mm}%
  \begin{center}{\headfont\bfseries\fontsize{14}{18}\selectfont #1\par}\end{center}\vspace{2pt}}
\newcommand\nuc[3]{\ensuremath{\prescript{#1}{#2}{\mathrm{#3}}}}

\begin{document}
\quiztitle{Year 12 Physics}{Short Answer Questions - Module 8 - From the Universe to the Atom:\\
  Marking Guidelines and Sample Answers}

Data used: $c = \qty{3.00e8}{m.s^{-1}}$, $h = \qty{6.626e-34}{J.s}$, $e = \qty{1.602e-19}{C}$,
$m_\mathrm{e} = \qty{9.109e-31}{kg}$, $m_\mathrm{p} = \qty{1.673e-27}{kg}$,
$\qty{1}{u} = \qty{1.661e-27}{kg} = \qty{931.5}{MeV}/c^2$, $R = \qty{1.097e7}{m^{-1}}$, $g = \qty{9.8}{m.s^{-2}}$.
In calculations, a correct answer from a correct method with an earlier error carried forward still earns
the later marks.

% =====================================================================
\anstopic{Topic 1 - Origins of the Elements}

% ---- Q1
\ansq
\anspart{a}
\begin{crit}
  \critrow{Correctly calculates the gradient of the line of best fit}{2}
  \critrow{Identifies $H_0$ as the gradient, OR calculates $\frac{v}{d}$ for one data point that is not on the line}{1}
\end{crit}
\sample
$v = H_0 d$, so $H_0$ is the gradient of the line:
$H_0 = \dfrac{6800 - 0}{100 - 0} = \qty{68}{km.s^{-1}.Mpc^{-1}}$ (accept 66 to 70).
\anspart{b}
\begin{crit}
  \critrow{Converts $H_0$ to SI units AND correctly calculates the age in years}{2}
  \critrow{Shows one correct step}{1}
\end{crit}
\sample
$H_0 = \dfrac{\qty{68e3}{m.s^{-1}}}{\qty{3.09e22}{m}} = \qty{2.20e-18}{s^{-1}}$. \quad
$t = \dfrac{1}{H_0} = \qty{4.5e17}{s} = \dfrac{\num{4.5e17}}{\num{3.16e7}} = \qty{1.4e10}{years}$.
\anspart{c}
\begin{crit}
  \critrow{Links recession proportional to distance to an expanding universe that began from a hot,
    dense state}{1}
\end{crit}
\sample
Galaxies are moving away from us with speeds proportional to their distance, so the universe is
expanding. Running the expansion backwards, all the matter was once together in an extremely hot, dense
state, as the Big Bang theory says.

% ---- Q2
\ansq
\begin{crit}
  \critrow{Describes the sequence: energy forms matter ($E = mc^2$) as the universe expands and cools;
    quarks and leptons; protons and neutrons; light nuclei; neutral atoms}{4}
  \critrow{Describes most of the sequence}{3}
  \critrow{Describes some stages, or the conversion of energy into matter}{2}
  \critrow{Gives one relevant point}{1}
\end{crit}
\sample
\begin{itemize}
  \item At first the universe was extremely hot and dense, and filled with radiation. As it expanded, it
        cooled. Energy was converted into particle--antiparticle pairs ($E = mc^2$); a small excess of
        matter over antimatter remained.
  \item In the first fractions of a second, the universe contained quarks, leptons (such as electrons)
        and photons.
  \item After about \qty{e-6}{s}, it was cool enough for quarks to combine into protons and neutrons.
  \item In the first few minutes, protons and neutrons fused into light nuclei: mainly helium-4, with a
        little deuterium and lithium. About 75\% of the mass stayed as hydrogen nuclei (protons).
  \item After about \num{380000} years, the universe was cool enough (about \qty{3000}{K}) for nuclei to
        capture electrons and form neutral atoms. Light could then travel freely; this light is now the
        cosmic microwave background.
\end{itemize}

% ---- Q3
\ansq
\anspart{a}
\begin{crit}
  \critrow{Identifies S, with a justification based on its position on the main sequence}{1}
\end{crit}
\sample
Star S. It is at the top left of the main sequence (the hottest and most luminous main-sequence star), and
along the main sequence, luminosity and temperature increase with mass.
\anspart{b}
\begin{crit}
  \critrow{Compares both temperature AND radius, with reasons from the diagram}{2}
  \critrow{Compares temperature OR radius}{1}
\end{crit}
\sample
Star P is much hotter than star Q (it is further left). But P is much less luminous than Q, even though
each square metre of P's surface emits more energy. So P must be very small (a white dwarf), and Q, which
is cool but very luminous, must be very large (a red giant). Q has a much larger radius than P.
\anspart{c}
\begin{crit}
  \critrow{Describes movement to the red-giant region (up and to the right) AND then to the white-dwarf
    region (down and to the left)}{2}
  \critrow{Describes one of these stages}{1}
\end{crit}
\sample
When the hydrogen in its core runs out, the Sun's core contracts and its outer layers expand and cool. It
becomes a red giant: its point moves up and to the right (more luminous, cooler). Later, after it has
fused helium and thrown off its outer layers (a planetary nebula), the hot core is left as a white dwarf:
its point moves to the bottom left (hot but very faint), and it then slowly cools and fades.

% ---- Q4
\ansq
\anspart{a}
\begin{crit}
  \critrow{Correctly calculates the energy in MeV}{2}
  \critrow{Correctly calculates the mass defect}{1}
\end{crit}
\sample
$\Delta m = 4(1.007825) - 4.002603 = \qty{0.028697}{u}$. \quad
$E = 0.028697 \times 931.5 = \qty{26.7}{MeV}$.
\anspart{b}
\begin{crit}
  \critrow{Gives a correct similarity AND a correct difference}{2}
  \critrow{Gives a similarity OR a difference}{1}
\end{crit}
\sample
Similar: both convert four hydrogen nuclei into one helium-4 nucleus (with positrons and neutrinos),
releasing the same energy (about \qty{26.7}{MeV}). Different: the proton--proton chain fuses protons
directly and dominates in stars like the Sun and smaller. The CNO cycle uses carbon, nitrogen and oxygen
nuclei as catalysts (they are not used up). Because these nuclei have larger charges, it needs a higher
core temperature, so it dominates in main-sequence stars more than about 1.3 times the mass of the Sun.

% =====================================================================
\anstopic{Topic 2 - Structure of the Atom}

% ---- Q5
\ansq
\anspart{a}
\begin{crit}
  \critrow{Correctly calculates $E$ AND the speed}{2}
  \critrow{Correctly calculates $E$, OR uses $v = \frac{E}{B}$}{1}
\end{crit}
\sample
$E = \dfrac{V}{d} = \dfrac{400}{0.020} = \qty{2.0e4}{V.m^{-1}}$. Undeflected: $qE = qvB$, so
$v = \dfrac{E}{B} = \dfrac{\num{2.0e4}}{\num{8.0e-4}} = \qty{2.5e7}{m.s^{-1}}$.
\anspart{b}
\begin{crit}
  \critrow{Correctly calculates the charge-to-mass ratio}{2}
  \critrow{Uses $qvB = \frac{mv^2}{r}$}{1}
\end{crit}
\sample
$qvB = \dfrac{mv^2}{r}$, so $\dfrac{q}{m} = \dfrac{v}{Br} = \dfrac{\num{2.5e7}}{(\num{8.0e-4})(0.18)}
= \qty{1.7e11}{C.kg^{-1}}$.
\anspart{c}
\begin{crit}
  \critrow{States that the particles' charge-to-mass ratio is about 1800 times larger, so (for a similar
    charge) they are much lighter than atoms: subatomic particles found in all atoms}{1}
\end{crit}
\sample
The ratio is about 1800 times that of a hydrogen ion. If the charges are similar in size, the cathode-ray
particles have about $\frac{1}{1800}$ of the mass of the lightest atom. So they are subatomic particles
(electrons), and since the result was the same for every cathode material, they are found in all atoms.

% ---- Q6
\ansq
\anspart{a}
\begin{crit}
  \critrow{Correctly calculates the charge AND states that it is negative, with a reason}{2}
  \critrow{Correctly calculates the size of the charge, OR states the sign with a reason}{1}
\end{crit}
\sample
Stationary: $qE = mg$, so $q = \dfrac{mg}{E} = \dfrac{(\num{3.27e-15})(9.8)}{\num{4.0e4}} = \qty{8.0e-19}{C}$.
The electric force must be upwards, opposite to the downward field, so the charge is negative.
\anspart{b}
\begin{crit}
  \critrow{5 electrons}{1}
\end{crit}
\sample
$n = \dfrac{\num{8.0e-19}}{\num{1.602e-19}} = 5$.
\anspart{c}
\begin{crit}
  \critrow{States that every measured charge was a whole-number multiple of one smallest value, $e$}{1}
\end{crit}
\sample
Every charge he measured was a whole-number multiple of about \qty{1.6e-19}{C}, never a fraction of it.
So charge comes in units of $e$, the charge of one electron.

% ---- Q7
\ansq
\begin{crit}
  \critrow{Describes the experiment AND the results, and explains how each result led to the nuclear
    model, in contrast with the plum-pudding model}{5}
  \critrow{Describes the experiment and results, and links them to the nuclear model}{4}
  \critrow{Describes the experiment and results, with a limited link to the model}{3}
  \critrow{Describes the experiment or the results}{2}
  \critrow{Gives one relevant point}{1}
\end{crit}
\sample
\begin{itemize}
  \item \textbf{Experiment:} alpha particles from a radioactive source were fired at a very thin gold foil
        in a vacuum. A zinc sulfide screen, viewed through a microscope, was moved around the foil to count
        the flashes (scintillations) made by the alpha particles at different angles.
  \item \textbf{Expected (plum-pudding model):} the positive charge is spread out, so its field is weak,
        and all the alpha particles should pass through with only tiny deflections.
  \item \textbf{Results:} most alpha particles went straight through or were deflected only slightly, but
        a very small fraction (roughly 1 in 10\,000) were deflected through more than \ang{90}, some almost
        straight back.
  \item \textbf{Rutherford's conclusions:} since most passed straight through, the atom is mostly empty
        space. To turn an alpha particle back needs a very strong force, so all the positive charge and
        almost all the mass must be concentrated in a tiny, dense nucleus (about $10^{-14}$ m across,
        compared with about $10^{-10}$ m for the atom). The rare large deflections show the nucleus is
        tiny. The electrons orbit the nucleus at a large distance, like planets around the Sun.
\end{itemize}

% ---- Q8
\ansq
\begin{crit}
  \critrow{Explains that the radiation is neutral, that gamma rays could not give the protons their
    observed energy, AND that conservation of momentum and energy in collisions showed a neutral particle
    with about the mass of a proton}{4}
  \critrow{Explains most of these points}{3}
  \critrow{Describes the experiment with some link to the neutron}{2}
  \critrow{Gives one relevant point}{1}
\end{crit}
\sample
Alpha particles hitting beryllium produced a very penetrating radiation that was not deflected by
electric or magnetic fields, so it was neutral. Others thought it was gamma radiation. When it hit
paraffin wax (rich in hydrogen), it knocked out protons with energies of about \qty{5.7}{MeV}. Using
conservation of energy and momentum, a gamma photon would need about \qty{50}{MeV} to do this, far more
than the reaction could supply. Chadwick also measured the recoil of nitrogen nuclei. Applying
conservation of momentum and energy to both collisions (like billiard balls), he showed that the radiation
was a neutral particle with a mass almost equal to the proton's: the neutron.
($\nuc{9}{4}{Be} + \nuc{4}{2}{He} \rightarrow \nuc{12}{6}{C} + \nuc{1}{0}{n}$)

% =====================================================================
\anstopic{Topic 3 - Quantum Mechanical Nature of the Atom}

% ---- Q9
\ansq
\anspart{a}
\begin{crit}
  \critrow{Correctly calculates the wavelength}{2}
  \critrow{Substitutes correctly into the Rydberg equation}{1}
\end{crit}
\sample
$\dfrac{1}{\lambda} = R\left(\dfrac{1}{2^2} - \dfrac{1}{4^2}\right) = (\num{1.097e7})(0.1875) = \qty{2.057e6}{m^{-1}}$,
so $\lambda = \qty{4.86e-7}{m}$ (\qty{486}{nm}).
\anspart{b}
\begin{crit}
  \critrow{Calculates the photon energy in eV AND compares it with the energy-level difference}{2}
  \critrow{Calculates the photon energy, OR the energy-level difference}{1}
\end{crit}
\sample
$E = \dfrac{hc}{\lambda} = \dfrac{(\num{6.626e-34})(\num{3.00e8})}{\num{4.86e-7}} = \qty{4.09e-19}{J}
= \qty{2.55}{eV}$. From the diagram: $E_4 - E_2 = -0.85 - (-3.40) = \qty{2.55}{eV}$. They agree.
\anspart{c}
\begin{crit}
  \critrow{Visible light}{1}
\end{crit}
\sample
Visible light (blue-green).

% ---- Q10
\ansq
\begin{crit}
  \critrow{Describes the strengths AND limitations of Bohr's model, and makes a judgement}{4}
  \critrow{Describes strengths and limitations}{3}
  \critrow{Describes a strength or a limitation}{2}
  \critrow{Gives one relevant point}{1}
\end{crit}
\sample
\begin{itemize}
  \item \textbf{Strengths:} Bohr combined Rutherford's nuclear model with Planck's quantum idea. His
        postulated stationary states, in which electrons do not radiate, accounted for the stability of atoms.
        Transitions between energy levels ($hf = E_i - E_f$) explained the line spectrum of hydrogen, and
        his model calculated the hydrogen wavelengths accurately (it gave the Rydberg equation).
  \item \textbf{Limitations:} it worked only for hydrogen and other one-electron ions. It could not
        explain the spectra of atoms with more electrons, the different intensities of the lines, or the
        splitting of lines (fine structure, and the Zeeman effect in a magnetic field). Its postulates were
        not explained: it mixed classical orbits with quantum rules.
  \item \textbf{Judgement:} Bohr's model was a very important step, because it introduced quantised
        energy levels, which are still used today. But it was incomplete, and it was replaced by the
        quantum mechanical model of de Broglie and Schrödinger.
\end{itemize}

% ---- Q11
\ansq
\anspart{a}
\begin{crit}
  \critrow{Correctly calculates the speed}{2}
  \critrow{Uses $qV = \frac{1}{2}mv^2$}{1}
\end{crit}
\sample
$qV = \frac{1}{2}mv^2$, so $v = \sqrt{\dfrac{2qV}{m}} = \sqrt{\dfrac{2(\num{1.602e-19})(150)}{\num{9.109e-31}}}
= \qty{7.3e6}{m.s^{-1}}$.
\anspart{b}
\begin{crit}
  \critrow{Correctly calculates the wavelength}{2}
  \critrow{Uses $\lambda = \frac{h}{mv}$}{1}
\end{crit}
\sample
$\lambda = \dfrac{h}{mv} = \dfrac{\num{6.626e-34}}{(\num{9.109e-31})(\num{7.26e6})} = \qty{1.0e-10}{m}$.
\anspart{c}
\begin{crit}
  \critrow{States that the wavelength is similar to the spacing of atoms in a crystal}{1}
\end{crit}
\sample
Diffraction is significant when the wavelength is similar to the size of the gaps. The spacing between
atoms in a crystal is about \qty{e-10}{m}, the same as the electrons' wavelength, so the crystal acts like
a diffraction grating.

% ---- Q12
\ansq
\begin{crit}
  \critrow{Explains electron standing waves with a whole number of wavelengths around the orbit, links this
    to Bohr's stationary states, AND outlines one piece of evidence}{4}
  \critrow{Explains most of these points}{3}
  \critrow{Describes matter waves with some link to Bohr's model, OR describes the evidence}{2}
  \critrow{Gives one relevant point}{1}
\end{crit}
\sample
De Broglie proposed that all matter has a wavelength $\lambda = \frac{h}{mv}$. He treated an electron in an
orbit as a standing wave. A standing wave fits only if a whole number of wavelengths fits around the
circumference: $2\pi r = n\lambda$. Other orbits would make the wave interfere destructively with itself.
Substituting $\lambda = \frac{h}{mv}$ gives $mvr = \frac{nh}{2\pi}$, exactly Bohr's condition. Because a
standing wave does not change with time, the electron does not radiate, which explains why these states
are stationary.

\textbf{Evidence:} Davisson and Germer (1927) fired electrons at a nickel crystal and found strong
reflection at particular angles: a diffraction pattern. The wavelength calculated from the pattern
agreed with $\lambda = \frac{h}{mv}$. (Also accepted: G.\,P.\ Thomson's electron diffraction rings through
thin metal foil.)

% =====================================================================
\anstopic{Topic 4 - Properties of the Nucleus}

% ---- Q13
\ansq
\anspart{a}
\begin{crit}
  \critrow{6.0 days (accept 5.5 to 6.5 days)}{1}
\end{crit}
\sample
The curve falls from \qty{1200}{Bq} to \qty{600}{Bq} in 6.0 days, so the half-life is 6.0 days.
\anspart{b}
\begin{crit}
  \critrow{Correctly calculates the decay constant}{1}
\end{crit}
\sample
$\lambda = \dfrac{\ln 2}{t_{1/2}} = \dfrac{0.693}{6.0} = \qty{0.116}{day^{-1}}$ (or
$\dfrac{0.693}{6.0 \times 86\,400} = \qty{1.34e-6}{s^{-1}}$).
\anspart{c}
\begin{crit}
  \critrow{Correctly calculates the activity}{2}
  \critrow{Substitutes correctly into $A = A_0e^{-\lambda t}$}{1}
\end{crit}
\sample
$A = A_0 e^{-\lambda t} = 1200\,e^{-0.1155 \times 40} = \qty{12}{Bq}$. (Accept 10 to \qty{12}{Bq}:
rounding $\lambda$ to 0.12 gives about \qty{10}{Bq}.)
\anspart{d}
\begin{crit}
  \critrow{Explains that radioactive decay is random, so the number of decays in each time interval
    varies}{1}
\end{crit}
\sample
Radioactive decay is random. We cannot predict when a particular nucleus will decay, so the number of
decays measured in each time interval varies around the average value given by the curve. (Background
radiation and measurement uncertainty add to the scatter.)

% ---- Q14
\ansq
\begin{crit}
  \critrow{Explains that both fission and fusion produce nuclei with a higher binding energy per nucleon
    (closer to the peak near iron-56), so the total binding energy increases and the mass decreases,
    releasing energy ($E = \Delta m c^2$); refers to the graph}{4}
  \critrow{Explains one process fully and the other in part}{3}
  \critrow{Explains one process, OR identifies the peak and links it to stability}{2}
  \critrow{Gives one relevant point}{1}
\end{crit}
\sample
The binding energy per nucleon shows how tightly the nucleons are held. It is greatest near iron-56
(about \qty{8.8}{MeV}), the most stable nuclei.
\begin{itemize}
  \item \textbf{Fission:} uranium-235 has about \qty{7.6}{MeV} per nucleon. When it splits into two
        medium nuclei (around $A = 90$ and 140), they have about \qty{8.5}{MeV} per nucleon, closer to the
        peak. For about 235 nucleons, this increase is roughly $235 \times \qty{0.9}{MeV} \approx
        \qty{200}{MeV}$.
  \item \textbf{Fusion:} light nuclei such as hydrogen-2 (about \qty{1.1}{MeV} per nucleon) fuse to form
        helium-4 (about \qty{7.1}{MeV} per nucleon), a large increase.
\end{itemize}
In both cases, the products are more tightly bound than the reactants, so their total binding energy is
larger. This means the products have less mass, and the lost mass is released as energy ($E = \Delta m c^2$).

% ---- Q15
\ansq
\anspart{a}
\begin{crit}
  \critrow{$x = 3$}{1}
\end{crit}
\sample
Mass numbers: $235 + 1 = 141 + 92 + x$, so $x = 3$. (Atomic numbers: $92 = 56 + 36$.)
\anspart{b}
\begin{crit}
  \critrow{Correctly calculates the energy in MeV}{2}
  \critrow{Correctly calculates the mass defect}{1}
\end{crit}
\sample
$\Delta m = (235.0439 + 1.008665) - (140.9144 + 91.9262 + 3 \times 1.008665) = 236.0526 - 235.8666
= \qty{0.1860}{u}$. \quad $E = 0.1860 \times 931.5 = \qty{173}{MeV}$.
\anspart{c}
\begin{crit}
  \critrow{Correctly calculates the total energy}{2}
  \critrow{Correctly calculates the number of nuclei, OR the energy of one fission in joules}{1}
\end{crit}
\sample
Number of nuclei: $N = \dfrac{1.0}{235.0439 \times \num{1.661e-27}} = \num{2.56e24}$. \quad
Energy per fission: $173.2 \times \num{1.602e-13} = \qty{2.77e-11}{J}$. \quad
Total: $(\num{2.56e24})(\num{2.77e-11}) = \qty{7.1e13}{J}$.

% ---- Q16
\ansq
\begin{crit}
  \critrow{Compares all three types for nature, charge AND penetrating ability, and explains how they are
    deflected differently in a magnetic field (direction and amount)}{4}
  \critrow{Compares the three types and describes the deflection, with minor errors}{3}
  \critrow{Compares some properties, OR describes the deflection}{2}
  \critrow{Gives one relevant point}{1}
\end{crit}
\sample
\begin{center}
  \renewcommand\arraystretch{1.4}
  \begin{tabular}{|l|p{36mm}|p{30mm}|p{56mm}|}
    \hline
    & \textbf{Nature} & \textbf{Charge} & \textbf{Penetrating ability}\\ \hline
    Alpha & helium nucleus (2 protons, 2 neutrons) & $+2e$ & low: stopped by paper or a few cm of air\\ \hline
    Beta & fast electron ($\beta^-$) or positron ($\beta^+$) & $-e$ (or $+e$) & medium: stopped by a few mm of
    aluminium\\ \hline
    Gamma & high-energy EM photon & 0 & high: reduced by thick lead or concrete\\ \hline
  \end{tabular}
\end{center}
In a magnetic field at right angles to their motion, the charged particles follow curved paths
($r = \frac{mv}{qB}$). Alpha and beta-minus particles curve in opposite directions because their charges
have opposite signs. Beta particles curve much more sharply (smaller radius), because they have much less
momentum for their charge ($r = \frac{p}{qB}$; their charge-to-mass ratio is far larger). Gamma rays are
uncharged, so they are not deflected.

% =====================================================================
\anstopic{Topic 5 - Deep Inside the Atom}

% ---- Q17
\ansq
\anspart{a}
\begin{crit}
  \critrow{Shows that the charges add to $+e$}{1}
\end{crit}
\sample
$\left(+\frac{2}{3}\right) + \left(+\frac{2}{3}\right) + \left(-\frac{1}{3}\right) = +1$, so the charge is
$+e$.
\anspart{b}
\begin{crit}
  \critrow{$+e$}{1}
\end{crit}
\sample
Up quark: $+\frac{2}{3}e$; anti-down quark: $+\frac{1}{3}e$ (opposite to the down quark). Total: $+e$.
\anspart{c}
\begin{crit}
  \critrow{Classifies all four correctly}{2}
  \critrow{Classifies two or three correctly}{1}
\end{crit}
\sample
Electron: lepton. \quad Neutron: baryon (three quarks, udd). \quad $\pi^+$: meson (quark and antiquark).
\quad Photon: gauge boson (carries the electromagnetic force).

% ---- Q18
\ansq
\begin{crit}
  \critrow{Explains that a two-body decay would give electrons one fixed energy, that the continuous
    spectrum suggested missing energy (and momentum), that Pauli proposed a light neutral particle to
    conserve them, AND describes the quark-level decay (d $\rightarrow$ u; the W boson is not required)}{5}
  \critrow{Explains most of these points}{4}
  \critrow{Explains the neutrino prediction OR the quark-level decay well}{3}
  \critrow{Describes some relevant points}{2}
  \critrow{Gives one relevant point}{1}
\end{crit}
\sample
\begin{itemize}
  \item If a nucleus emitted only an electron (so there were just two products: the daughter nucleus and
        the electron), conservation of energy and momentum would give every electron the same kinetic
        energy. Instead, the electrons have a
        continuous range of energies, up to a maximum. It looked as if energy (and momentum) were not
        conserved.
  \item In 1930, Pauli proposed that a third particle is also emitted, carrying away the missing energy
        and momentum. It had to be neutral (charge was already conserved) and have very little or no mass
        (the maximum electron energy equals the total energy released). Fermi named it the neutrino. It
        interacts so weakly that it was not detected until 1956.
  \item At the quark level, one down quark in the neutron (udd) changes into an up quark, making a proton
        (uud). This happens through the weak force: the down quark emits a $\mathrm{W}^-$ boson, which
        quickly decays into an electron and an electron antineutrino:
        $\mathrm{n} \rightarrow \mathrm{p} + \mathrm{e}^- + \bar{\nu}_\mathrm{e}$. Charge and lepton number
        are conserved.
\end{itemize}

% ---- Q19
\ansq
\anspart{a}
\begin{crit}
  \critrow{Correctly calculates $\gamma$}{2}
  \critrow{Converts the energy to joules, OR calculates $mc^2$}{1}
\end{crit}
\sample
$E = (\num{6.8e12})(\num{1.602e-19}) = \qty{1.09e-6}{J}$. \quad
$mc^2 = (\num{1.673e-27})(\num{3.00e8})^2 = \qty{1.51e-10}{J}$. \quad
$\gamma = \dfrac{E}{mc^2} = \dfrac{\num{1.09e-6}}{\num{1.51e-10}} = \num{7.2e3}$.
\anspart{b}
\begin{crit}
  \critrow{Explains that collision energy is converted into the mass of new heavy particles ($E = mc^2$),
    AND/OR that high-energy particles have short wavelengths to probe small distances, with clear
    reasoning}{2}
  \critrow{States one reason without explanation}{1}
\end{crit}
\sample
In a collision, energy can be converted into mass ($E = mc^2$). Many particles are very massive (for
example, the Higgs boson, discovered at the LHC in 2012, has about 133 times the mass of a proton), so very
large collision energies are needed to create them. Also, high-energy particles have a very short de Broglie wavelength
($\lambda = \frac{h}{p}$), so they can probe very small distances, such as the structure inside protons.
\anspart{c}
\begin{crit}
  \critrow{Explains that a large $\gamma$ means $v$ is very close to $c$, but reaching $c$ would need
    infinite energy}{1}
\end{crit}
\sample
$\gamma = \num{7.2e3}$ means $\sqrt{1 - \frac{v^2}{c^2}}$ is tiny, so $v$ is extremely close to $c$
(about $0.99999999c$). But $\gamma$ becomes infinite at $v = c$, so reaching $c$ would need infinite
energy. Extra energy increases $\gamma$ a lot, but barely changes the speed.

% ---- Q20
\ansq
\begin{crit}
  \critrow{Analyses at least TWO pieces of evidence, explaining how each shows that protons and neutrons
    have internal structure (quarks)}{4}
  \critrow{Describes two pieces of evidence with some explanation}{3}
  \critrow{Describes one piece of evidence}{2}
  \critrow{Gives one relevant point}{1}
\end{crit}
\sample
\begin{itemize}
  \item \textbf{Deep inelastic scattering (Stanford, 1968):} high-energy electrons, with de Broglie
        wavelengths much smaller than a proton, were fired at protons. Many more were scattered through
        large angles than expected for a smooth, spread-out charge. Just as in the Geiger--Marsden
        experiment, this showed that the charge is concentrated in small, point-like particles inside the
        proton. Later experiments (early 1970s) showed that these have fractional charges: quarks.
  \item \textbf{The ``particle zoo'':} accelerators produced hundreds of new hadrons. Gell-Mann's
        ``Eightfold Way'' arranged them in patterns and predicted a missing particle, the $\Omega^-$, which
        was found in 1964. That year Gell-Mann and Zweig explained these patterns: every hadron is a
        combination of a few quarks (the $\Omega^-$ is sss).
  \item \textbf{Other evidence:} the neutron has no charge but has a magnetic moment, so it must contain
        moving charges. Beta decay, in which a neutron changes into a proton, is also consistent with the
        quark picture: one down quark changes into an up quark.
\end{itemize}
Together, this evidence shows that protons and neutrons are made of quarks, which are fundamental, even
though free quarks are never seen on their own.

\end{document}
